% 原书第135—138页：导言、7.1.1—7.1.2。
\chapter{学习多元因果模型}
\label{ch:learning-multivariate}
\newcounter{chseventransnote}
\providecommand{\chseventnote}[1]{\stepcounter{chseventransnote}\begingroup
\renewcommand{\thefootnote}{译注\arabic{chseventransnote}}\footnote{#1}\addtocounter{footnote}{-1}\endgroup}

与第4章一样，下面转向因果模型的学习问题。首先，第7.1节讨论在什么假设下，能够从联合分布恢复全部或部分图结构，即“结构可识别性”。其中一些结果可以从前面讨论的二元情形推广而来。与二元情形一样，目前尚未完整刻画能够保证可识别性的全部假设，未来的研究也许会揭示有前景的替代方案。接着，第7.2节介绍根据有限数据集估计图的方法和算法，例如基于独立性的方法及基于得分的方法，即“结构识别”。

与二元情形一样，这里仍面临\textbf{结构因果模型（structural causal model，SCM）}类过于灵活的问题。给定随机变量 $X=(X_1,\ldots,X_d)$ 的分布 $P_X$，不同的结构因果模型是否可能导出这一分布？下面的命题回答了这个问题：实际上，对于许多不同的图结构，通常都存在能够诱导出分布 $P_X$ 的结构因果模型。\footnote{类似于命题7.1的论断，可见 Druzdzel and Simon［1993］及 Druzdzel and van Leijen［2001］。}

\begin{proposition}[图结构的非唯一性]
考虑随机向量 $X=(X_1,\ldots,X_d)$，其分布 $P_X$ 关于 Lebesgue 测度具有密度。假设 $P_X$ 关于图 $\mathcal G$ 满足马尔可夫性。那么，存在一个图为 $\mathcal G$ 的结构因果模型 $\mathfrak C=(\mathcal S,P_N)$，能够导出分布 $P_X$。
\end{proposition}
\begin{proof}
见附录C.9。
\end{proof}

特别地，给定任意一个完全\textbf{有向无环图（directed acyclic graph，DAG）}，都能找到一个相应的结构因果模型，导出手头的分布。因此，与二元情形一样，要获得可识别性结果，显然还需要进一步的假设。下一节将讨论其中一些假设。

\section{结构可识别性}
\subsection{忠实性}

如果分布 $P_X$ 关于底层有向无环图 $\mathcal G^0$ 满足马尔可夫性和\textbf{忠实性（faithfulness）}，那么，图 $\mathcal G^0$ 中的 d-分离陈述，与分布中相应的条件独立陈述之间，就具有一一对应关系。因此，可以排除所有不属于 $\mathcal G^0$ 正确马尔可夫等价类的图，因为这些图所施加的 d-分离集合，与 $P_X$ 中的条件独立集合不相等。马尔可夫条件与忠实性只限制联合分布中的条件独立关系，因此，同样显然的是，我们无法区分两个马尔可夫等价的图，也就是导出完全相同条件独立集合的两个图，例如见原书第103页的图6.4。总之，在马尔可夫条件与忠实性成立时，$\mathcal G^0$ 的马尔可夫等价类可以从 $P_X$ 中识别，并由 $\mathrm{CPDAG}(\mathcal G^0)$ 表示［例如 Spirtes et al., 2000］。

\begin{lemma}[马尔可夫等价类的可识别性]
假设 $P_X$ 关于 $\mathcal G^0$ 满足马尔可夫性和忠实性。那么，对于每一个图 $\mathcal G\in\mathrm{CPDAG}(\mathcal G^0)$，都能找到一个导出分布 $P_X$ 的结构因果模型。此外，不存在 $\mathcal G\notin\mathrm{CPDAG}(\mathcal G^0)$，使得 $P_X$ 关于 $\mathcal G$ 同时满足马尔可夫性和忠实性。
\end{lemma}
\begin{proof}
第一条结论直接由命题7.1推出；第二条结论由定义6.24中的马尔可夫等价性定义得到。
\end{proof}

\textbf{基于独立性的方法（independence-based methods）}，也称\textbf{基于约束的方法（constraint-based methods）}，假设分布关于底层图满足马尔可夫性和忠实性，再估计正确的马尔可夫等价类，见第7.2.1节。

例6.42已经说明，对于高斯分布，因果效应可以由式 (6.20) 中的一个数概括。如果只知道图的马尔可夫等价类，而不知道正确的图，那么这一量就不再具有可识别性。不过，仍然可以给出它的界［Maathuis et al., 2009］。

\subsection{加性噪声模型}

命题7.1表明，一个给定的分布，可以由具有不同图结构的多个结构因果模型导出。不过，对于其中许多图结构，结构赋值中的函数 $f_j$ 会相当复杂。事实表明，如果不允许函数任意复杂，也就是限制函数类，就能够获得非平凡的可识别性结果。与第4章中一样，在后面的第7.1.4和7.1.5节中，我们假定噪声以加性方式作用。

\begin{definition}[加性噪声模型]
如果结构因果模型 $\mathfrak C$ 的结构赋值具有如下形式，
\begin{equation}
 X_j:=f_j(\PA_j)+N_j,\qquad j=1,\ldots,d,
 \tag{7.1}
\end{equation}
即噪声是加性的，就称其为\textbf{加性噪声模型（additive noise model，ANM）}。为简单起见，进一步假定函数 $f_j$ 可微，且噪声变量 $N_j$ 具有严格为正的密度。\footnote{例如，这两个条件能够保证 $X_1,\ldots,X_d$ 的联合分布也具有严格为正的密度。}
\end{definition}

下面一些可识别性结果需要假设\textbf{因果极小性（causal minimality）}，见定义6.33。对于加性噪声模型，这意味着：每个函数 $f_j$ 对其任何一个自变量都不是常函数。直观地说，函数应当真正“依赖”其自变量。下面命题的证明见附录C.10。

\begin{proposition}[因果极小性与加性噪声模型]
考虑由模型 (7.1) 诱导的分布，并假定函数 $f_j$ 对其任何一个自变量都不是常函数。也就是说，对于所有 $j$ 和 $i\in\PA_j$，都存在变量 $\PA_j\setminus\{i\}$ 的某个取值 $\pa_{j,-i}$，以及某些 $x_i\neq x'_i$，使得
\[
 f_j(\pa_{j,-i},x_i)\neq f_j(\pa_{j,-i},x'_i).
\]
那么，联合分布关于相应图满足因果极小性。反之，如果存在节点 $j,i$，使得对所有 $\pa_{j,-i}$，函数 $f_j(\pa_{j,-i},\cdot)$ 都为常函数，那么因果极小性就不成立。
\end{proposition}

注6.6已论证，可以将讨论限于对任一自变量都不是常函数的函数，另见命题6.49。现在我们又看到，对于噪声具有全支撑的加性噪声模型，这一限制蕴含因果极小性。

在式 (7.1) 所描述的受限结构因果模型类中，是否能够得到完整的结构可识别性？答案仍然是否定的。例如，定理4.2和习题7.13表明，如果分布由一个线性高斯结构因果模型诱导，就不一定能够恢复正确的图。不过，从以下意义上说，这种情况是例外：对于几乎所有其他函数与分布的组合，都能获得可识别性。所有不可识别情形都已经得到刻画［Zhang and Hyvärinen, 2009；Peters et al., 2014］。另一个不同于线性高斯情形的不可识别例子，见图4.2右图；其细节可见 Peters et al.［2014，例25］。表7.1汇总了部分已知的可识别性结果。

\begin{table}[htbp]
\centering
\small
\setlength{\tabcolsep}{4pt}
\begin{tabular}{@{}llccl@{}}
\toprule
\multicolumn{2}{c}{结构赋值类型}&函数条件&DAG可识别&参见\\
\midrule
（一般）SCM&$X_j:=f_j(X_{\PA_j},N_j)$&—&否&命题7.1\\
ANM&$X_j:=f_j(X_{\PA_j})+N_j$&非线性&是&定理7.7（i）\\
CAM&$X_j:=\sum_{k\in\PA_j}f_{jk}(X_k)+N_j$&非线性&是&定理7.7（ii）\\
线性高斯&$X_j:=\sum_{k\in\PA_j}\beta_{jk}X_k+N_j$&线性&否&习题7.13\\
线性高斯、误差等方差&$X_j:=\sum_{k\in\PA_j}\beta_{jk}X_k+N_j$&线性&是&命题7.5\\
\bottomrule
\end{tabular}
\caption{高斯噪声下若干已知可识别性结果的汇总。对于非高斯噪声，也有可识别性结果，但其技术细节更复杂。}
\label{tab:ch07-identifiability}
\end{table}


再次说明，加性噪声模型框架有多种扩展。例如，Zhang and Hyvärinen［2009］允许对变量作后非线性变换，而 Peters et al.［2011a］考虑了离散变量的加性噪声模型。

一般而言，非线性加性噪声模型对边缘化不封闭。也就是说，即使 $P_{X,Y,Z}$ 容许从 $X$ 到 $Y$、从 $Y$ 到 $Z$ 的加性噪声模型，$P_{X,Z}$ 也不一定容许从 $X$ 到 $Z$ 的加性噪声模型。实际中，我们可能观测不到因果路径上的中间变量，因此，这一点可能限制加性噪声模型的适用范围。对于物理学实验，可以认为任何影响都是经由无穷多个中间变量传播的。因此，直接效应与间接效应并无绝对的定义，它们总是相对于所观测的变量集合而言。在这个意义上，加性噪声模型只能被看作良好的近似。

下面三个小节将更详细地讨论三个具有可识别性的具体例子：误差方差相等的线性高斯情形，第7.1.3节；线性非高斯情形，第7.1.4节；非线性高斯情形，第7.1.5节。虽然已经有更一般的结果［Peters et al., 2014］，这里仍集中讨论这些例子，因为它们的精确条件比较容易陈述。我们略去证明，重点介绍结论。大多数证明都可以建立在 Peters et al.［2011b］所发展的技术之上；这些技术使第4章介绍的许多二元可识别性结果能够推广到多元情形。
% 原书本段先列三个例子，随后写“those two examples”；中文使用“这些例子”避免数量自相矛盾。
